\section{Chapter Overview}

This example chapter demonstrates a common structure for a research study: motivation, definitions, a conceptual model, mathematical formulation, and a short summary. It also shows how to include an externally compiled TikZ figure from the chapter's \texttt{figs} folder.

\section{Conceptual Model}

Assume that a study evaluates several alternatives using three normalized criteria: quality, efficiency, and usability. Let \(q_i\), \(e_i\), and \(u_i\) represent the three scores for alternative \(i\). A simple weighted score is

\begin{equation}
\label{eq:weighted-score}
S_i = w_q q_i + w_e e_i + w_u u_i,
\end{equation}

where \(w_q+w_e+w_u=1\) and each weight is nonnegative. Equation~\eqref{eq:weighted-score} is not presented as a universal research model; it is an example of equation formatting, labeling, and cross-referencing.

The symbols can be added to the nomenclature:
\nomenclature{$S_i$}{Overall score for alternative $i$}
\nomenclature{$w_q,w_e,w_u$}{Weights assigned to the evaluation criteria}

\section{External TikZ Figure}

Figure~\ref{fig:research-process} is defined in \texttt{paper1/figs/Graph.tex}. Because that file uses the \texttt{standalone} document class, it can be compiled independently to inspect the diagram without compiling the full dissertation. The resulting figure is then included below.

\begin{figure}[ht]
    \centering
    \includestandalone[width=0.95\textwidth]{paper1/figs/Graph}
    \caption{A generic research process created in a standalone TikZ file.}
    \label{fig:research-process}
\end{figure}

\section{Example Table}

Table~\ref{tab:sample-results} contains invented values and demonstrates a caption, label, header, and numeric columns. Replace both the labels and values with actual evidence.

\begin{table}[ht]
\centering
\caption{Sample performance results using fictional data.}
\label{tab:sample-results}
\footnotesize
\resizebox{\columnwidth}{!}{%
\begin{tabular}{l|ccc}
\hline
\textbf{Method} & \textbf{Quality (\%)} & \textbf{Runtime (s)} & \textbf{Overall Score} \\
\hline
Baseline & 82.5 & 15.3 & 0.71 \\
Method A & 87.9 & 12.8 & 0.79 \\
Method B & 90.4 & 10.6 & 0.84 \\
Proposed & 93.1 & 9.4  & 0.89 \\
\hline
\end{tabular}}
\end{table}

The prose following a table should explain the important pattern rather than merely repeat every cell. For the fictional values in Table~\ref{tab:sample-results}, the proposed method has the highest quality and overall score as well as the shortest runtime. An actual dissertation should also report uncertainty, statistical tests, effect sizes, or qualitative evidence as appropriate.

\section{Chapter Summary}

This chapter illustrated equations, symbols, cross-references, a standalone TikZ figure, and a formatted table. Chapter~\ref{ch:paper2} extends the demonstration with algorithm and theorem environments.
